\documentclass[11pt,a4paper]{article}

% ── Packages ──────────────────────────────────────────────────────────────────
\usepackage[utf8]{inputenc}
\usepackage[T1]{fontenc}
\usepackage{amsmath,amssymb,amsfonts}
\usepackage{graphicx}
\usepackage{booktabs}
\usepackage{multirow}
\usepackage{hyperref}
\usepackage[margin=1in]{geometry}
\usepackage{natbib}
\usepackage{authblk}
\usepackage{caption}
\usepackage{subcaption}
\usepackage{xcolor}
\usepackage{algorithm}
\usepackage{algpseudocode}
\usepackage{enumitem}

\hypersetup{
  colorlinks=true,
  linkcolor=blue!70!black,
  citecolor=blue!70!black,
  urlcolor=blue!70!black
}

% ── Title ─────────────────────────────────────────────────────────────────────
\title{%
  Affect as Precision Augmentation for Social Inference:\\
  Per-Partner Metacognitive States in Multi-Agent Active Inference%
}

\author[1]{Harshil Shah}
\affil[1]{}

\date{March 2026}

% ══════════════════════════════════════════════════════════════════════════════
\begin{document}
\maketitle

% ── Abstract ──────────────────────────────────────────────────────────────────
\begin{abstract}
We present a computational model of affect as per-partner metacognitive precision tracking in multi-agent active inference. Each social partner is assigned a slow-timescale affective state $\beta_k$ that summarizes the running reliability of the agent's predictive model for that partner. This state augments policy evaluation by precision-weighting expected free energy estimates with partner-specific confidence information. In a volatile multi-partner trust game, we show that: (1)~affect provides \emph{orthogonal} augmentation beyond explicit planning depth---under sophisticated inference, non-affective planners from horizon $\tau = 2$ to $\tau = 8$ are statistically indistinguishable, yet adding the affective signal yields $d = 0.64$ improvement at every depth; (2)~ablating affect reproduces the somatic marker deficit---intact partner knowledge with impaired decision deployment; (3)~precision tracking and reward averaging occupy distinct informational niches, with precision tracking excelling under betrayal stress ($d = 0.59$) and reward averaging dominating in continuous-investment settings ($d = 1.72$); (4)~Bayesian model comparison decisively favors the affective model under volatility (log$_{10}$ BF $= 3.0$); (5)~augmentation generalizes across Prisoner's Dilemma, Stag Hunt, and Chicken games ($d > 1.0$ under betrayal in all three); and (6)~clinical parameter perturbations (alexithymia, borderline, depression) produce qualitatively distinct behavioral phenotypes across environments---in the Stag Hunt, borderline is the first phenotype with decisive Bayesian evidence of impairment ($d = 0.72$, log$_{10}$ BF $= -2.89$), while alexithymia is paradoxically protective. The mechanism occupies a distinct computational niche from both outside-in empathy and cognitive theory of mind: inside-out metacognitive monitoring of one's own social inference channel. Affect encodes information orthogonal to both planning depth and reward history: partner-specific model fitness that improves social decisions through a channel inaccessible to explicit deliberation alone.
\end{abstract}

\noindent\textbf{Keywords:} active inference, affect, social cognition, precision weighting, somatic markers, multi-agent systems, computational psychiatry

% ══════════════════════════════════════════════════════════════════════════════
\section{Introduction}
\label{sec:intro}

Social cognition presents a fundamental computational bottleneck. When an agent interacts with multiple social partners whose strategies may change unpredictably, optimal decision-making requires maintaining probabilistic models of each partner's hidden states, simulating forward through multiple rounds of interaction, and evaluating expected outcomes under candidate policies. The computational cost grows exponentially with planning horizon, yet humans navigate these environments fluently---suggesting they employ computational shortcuts that preserve decision quality without exhaustive forward planning.

Damasio's somatic marker hypothesis \citep{damasio1994descartes} proposes that emotional signals provide exactly this shortcut. Through the ventromedial prefrontal cortex (vmPFC) and its connections to the amygdala and insula, the brain generates rapid affective evaluations that bias decision-making before and below deliberate reasoning. The Iowa Gambling Task \citep{bechara1994insensitivity,bechara1997deciding} provides the key evidence: vmPFC-lesioned patients can verbally identify risky options but fail to develop the anticipatory somatic markers that healthy participants use to guide advantageous choices. The dissociation---intact knowledge but impaired deployment---is the behavioral signature our model aims to reproduce computationally in a social context.

Active inference \citep{friston2010free,parr2022active} provides the formal framework. Agents maintain generative models of observation causes and act to minimize variational free energy. Policy selection is governed by expected free energy (EFE), decomposing into pragmatic value (preference alignment) and epistemic value (uncertainty reduction). \citet{hesp2021deeply} formalized affect within this framework: valence encodes expected model precision (subjective fitness), with affective charge driven by signed prediction errors. Critically, their model operates in a single-agent context. What is absent is a per-partner metacognitive affective state that compresses interaction history into a precision signal for multi-agent social inference.

Existing multi-agent active inference work covers theory of mind as inference over others' generative models \citep{friston2015duet,yoshida2008game}, empathy as model coupling \citep{pitliya2025empathy,matsumura2024active}, and factorized multi-agent models \citep{ruizserra2025factorised}. These approaches define three distinct computational strategies for social cognition. \emph{Outside-in empathy} \citep{pitliya2025empathy} couples agents' generative models so that one agent's affective or perceptual states directly inform another's inference---the computational direction flows from partner to self. \emph{Cognitive theory of mind} \citep{yoshida2008game} maintains an explicit model of the partner's inference process, recursively simulating their beliefs and policies. The present work introduces a third strategy: \emph{inside-out metacognitive monitoring}, where the agent tracks the reliability of its own predictive model for each partner, without modeling the partner's internal states or coupling to their generative process. This triple dissociation---outside-in empathy, cognitive ToM, and inside-out precision monitoring---carves the space of social inference mechanisms at its computational joints. None of the existing approaches provides a per-partner metacognitive affective state that: (a)~compresses interaction history into a precision signal, (b)~provides orthogonal augmentation beyond what explicit planning depth recovers, and (c)~whose ablation reproduces the somatic marker deficit pattern.

\paragraph{Contributions.} We introduce three novel components:
\begin{enumerate}[nosep]
  \item \textbf{Per-partner metacognitive affective state in multi-agent active inference.} Not a global mood or shared precision parameter, but a learned state for each social partner compressing that specific interaction history into a precision signal.
  \item \textbf{Affect as orthogonal augmentation beyond planning depth.} Under sophisticated inference, explicit planning depth ($\tau = 2$ through $\tau = 8$) produces identical non-affective performance; the affective signal adds evaluative information that no amount of additional lookahead recovers.
  \item \textbf{Computational vmPFC lesion in a social context.} The ablation reproduces intact social knowledge with impaired deployment, extending the IGT literature to multi-agent settings.
\end{enumerate}

% ══════════════════════════════════════════════════════════════════════════════
\section{Model}
\label{sec:model}

\subsection{Three-Level Architecture}

The model has three hierarchical levels operating on different timescales:

\textbf{Level 1---Observations} (per-trial). Trial-by-trial interaction outcomes: partner's action, own payoff. Updated every round.

\textbf{Level 2---Cognitive social model} (parameter learning). Latent states encoding each partner's type/strategy (cooperator, reciprocator, exploiter, random). Dirichlet parameters accumulate evidence over many rounds.

\textbf{Level 3---Per-partner affective state} (metacognitive integration). A slow state $\beta_k$ for each partner $k$ summarizing the running quality of the agent's social model. Updated by an exponentially-weighted moving average of prediction error magnitudes, resistant to single-trial fluctuations.

\subsection{Generative Model}

The agent maintains a POMDP generative model for each social partner $k$:
\begin{equation}
  P(o_{1:T}, s_{1:T}, \pi) = P(\pi) \prod_{t=1}^{T} P(o_t \mid s_t)\, P(s_t \mid s_{t-1}, \pi)
\end{equation}
with hidden states factoring hierarchically: $s_t = (s_t^{(1)}, s_t^{(2)}, s_t^{(3)})$ for trial state, partner type, and affective state respectively.

State inference exploits the analytical solution to variational free energy minimization: for a categorical $q(s)$, the VFE-minimizing posterior equals the Bayesian posterior $q^*(s) \propto P(o \mid s) P(s)$. The implementation uses one-step Bayes with the likelihood matrix $\mathbf{A}$ and transition matrix $\mathbf{B}$.

\subsection{Expected Free Energy and Policy Selection}

For policy $\pi$ at future timestep $t$:
\begin{equation}
  G_t(\pi) = \underbrace{-\mathbb{E}_{Q(o_t|\pi)}[\ln P(o_t)]}_{\text{pragmatic value}} - \underbrace{\mathbb{E}_{Q(s_t|\pi)}[D_\text{KL}[Q(s_t|o_t,\pi) \| Q(s_t|\pi)]]}_{\text{epistemic value}}
\end{equation}

Policy selection follows a softmax: $P(\pi) = \sigma(-\gamma \cdot G(\pi))$, where $\gamma$ is the policy precision. The planner uses observation-branching sophisticated inference across all conditions, enumerating future binary observation sequences, updating beliefs pathwise via Bayes rule, and averaging pathwise EFE under path probabilities.

\subsection{Affective State Dynamics}

The affective state $\beta_k$ for partner $k$ is updated from a signed affective charge derived from prediction error:
\begin{equation}
  \beta_k^{(t+1)} = \lambda\, \beta_k^{(t)} + (1 - \lambda)\, \sigma\!\bigl(\phi(\epsilon_k^{(t)})\bigr)
\label{eq:beta_update}
\end{equation}
where $\lambda \in (0.5, 0.9)$ is a smoothing parameter, $\epsilon_k^{(t)} = 1 - P(o_t^{\mathrm{obs}} \mid s_k^{(2)})$ is the unsigned surprise from the level-2 model, and $\phi(\epsilon) = \alpha(\sigma_0^2 - \epsilon^2)$ converts surprise into affective charge. When actual squared surprise is below baseline $\sigma_0^2$, affect increases (model is accurate); when above, it decreases (model is failing).

This update extends \citeauthor{hesp2021deeply}'s variational treatment of precision dynamics to the per-partner social setting, building on the free-energy formalization of emotional valence by \citet{joffily2013emotional} and the interoceptive inference framework of \citet{seth2013interoceptive}. The surprise term is computed from the agent's own level-2 predictive beliefs, making the affective update belief-coupled rather than externally injected.

\subsection{Affect as Precision Weighting}

The affective agent augments standard EFE with a partner-specific precision weight:
\begin{equation}
  G(\pi) \approx \left(\sum_{t=1}^{\tau} G_t(\pi)\right) \cdot (1 + \mu\, \beta_k)
\label{eq:precision_weight}
\end{equation}
where $\mu$ scales the influence of affect on policy selection. The weight is keyed to the first partner implicated by the policy so the signal survives first-action marginalization.

\subsection{The vmPFC Lesion}

The computational lesion sets $\mu = 0$: the agent ``feels'' things about partners ($\beta_k$ updates normally) but feelings have no effect on action selection. This models a disconnection between metacognitive affective representation and decision circuitry, preserving intact level-2 posteriors while impairing behavioral deployment.

\subsection{Reward-Average Control (Condition 5)}

To distinguish precision tracking from cached value, a reward-average agent replaces $\beta_k$ with a simple exponential moving average of reward:
\begin{equation}
  \bar{R}_k^{(t+1)} = \lambda\, \bar{R}_k^{(t)} + (1 - \lambda)\, r_t, \qquad
  \tilde{R}_k = \tfrac{1}{2}(1 + \tanh(\bar{R}_k / \max |r|))
\end{equation}
The terminal signal $\tilde{R}_k$ enters the same precision-weighting mechanism as $\beta_k$, enabling controlled comparison of precision-based versus value-based augmentation.

\subsection{Discrete Variational Beta}

Phase 4 formalized $\beta_k$ as a discrete hidden state with $L = 5$ precision levels, explicit likelihood $P(\epsilon \mid \beta) \propto \exp(\alpha(\sigma_0^2 - \epsilon^2) \cdot \beta)$, and tridiagonal persistent transition dynamics. The posterior update uses the same categorical Bayesian machinery as partner-type inference:
\begin{equation}
  q(\beta_k^{(t)}) \propto P(\epsilon_k^{(t)} \mid \beta_k) \sum_{\beta_k^{(t-1)}} P(\beta_k^{(t)} \mid \beta_k^{(t-1)})\, q(\beta_k^{(t-1)})
\end{equation}
In stable conditions, continuous EMA and discrete formulations are indistinguishable ($d = 0.001$). Under betrayal, the discrete formulation underperforms moderately ($d = 0.41$) due to the transition matrix constraining single-step posterior shifts---a feature reflecting a stronger prior on precision stability.

% ══════════════════════════════════════════════════════════════════════════════
\section{Experimental Design}
\label{sec:design}

\subsection{Multi-Partner Trust Game}

The agent interacts with $K = 4$ partners over $N = 200$ rounds. Each round, a partner is selected (randomly or by agent choice), both choose cooperate or defect simultaneously, and payoffs follow a Prisoner's Dilemma matrix (mutual cooperation: +3/+3; sucker: $-1$/+5; temptation: +5/$-1$; mutual defection: +1/+1). Partners have latent types (cooperator, reciprocator, exploiter, random) that switch stochastically ($p_\text{switch} = 0.05$).

\subsection{Graded Investment Trust Game}

To test precision modulation in ambiguous decision spaces, we introduce a graded variant with six investment levels (0\%, 20\%, 40\%, 60\%, 80\%, 100\%) per partner, creating 24 total actions. This expands policy entropy from $< 0.01$ (binary) to $\sim 5.8$ (graded), activating the precision channel that is structurally inert under binary softmax saturation.

\subsection{Cross-Game Generalization}

We test three 2$\times$2 symmetric games with zero code changes---only the payoff matrix differs:

\begin{table}[ht]
\centering\small
\begin{tabular}{lccccc}
\toprule
Game & Mutual Coop & Sucker & Temptation & Mutual Defect & Character \\
\midrule
Prisoner's Dilemma & (3,3) & ($-1$,5) & (5,$-1$) & (1,1) & Defection-dominated \\
Stag Hunt & (5,5) & (0,2) & (2,0) & (2,2) & Coordination \\
Chicken & (3,3) & (1,5) & (5,1) & (0,0) & Anti-coordination \\
\bottomrule
\end{tabular}
\caption{Payoff matrices for cross-game generalization experiments.}
\label{tab:payoffs}
\end{table}

\subsection{Experimental Conditions}

\begin{table}[ht]
\centering\small
\begin{tabular}{clccc}
\toprule
ID & Condition & Horizon & Affect & Terminal signal \\
\midrule
C1 & Deep planner, no affect & $\tau=8$ & --- & --- \\
C2 & Affective shallow planner & $\tau=2$ & $\beta_k$ (precision) & $\beta_k$ \\
C3 & Lesioned affective & $\tau=2$ & $\beta_k$ (decoupled) & --- \\
C4 & Shallow planner, no affect & $\tau=2$ & --- & --- \\
C5 & Reward-average shallow & $\tau=2$ & --- & $\tilde{R}_k$ \\
C6 & Intermediate, no affect & $\tau=3$ & --- & --- \\
C7 & Intermediate, no affect & $\tau=4$ & --- & --- \\
C8 & Deep planner with affect & $\tau=8$ & $\beta_k$ (precision) & $\beta_k$ \\
C12 & Variational affective & $\tau=2$ & $q(\beta_k)$ (discrete) & $\hat{\beta}_k$ \\
\bottomrule
\end{tabular}
\caption{Experimental conditions. All use observation-branching sophisticated inference.}
\label{tab:conditions}
\end{table}

\subsection{Betrayal Stress Protocol}

To test robustness under volatility, we use an agent-chosen partner assignment with a scheduled cooperator$\to$exploiter switch at round 31. Stochastic type switches are disabled ($p_\text{switch} = 0$) to isolate the scheduled betrayal. This creates a regime where recent reward and prediction error dissociate---the critical test for precision tracking.

\subsection{Clinical Parameter Regimes}

Three clinical phenotypes are modeled as parameter perturbations, following the computational psychiatry framework of \citet{smith2021bayesian}:
\begin{itemize}[nosep]
  \item \textbf{Alexithymia} (C9): $\alpha = 0.1$ (blunted affective charge)
  \item \textbf{Borderline} (C10): $\alpha = 12$, $\lambda = 0.5$ (volatile, poorly smoothed precision)
  \item \textbf{Depression} (C11): $\beta_0 = 0.2$ (pessimistic initial precision prior)
\end{itemize}

% ══════════════════════════════════════════════════════════════════════════════
\section{Results}
\label{sec:results}

\subsection{Orthogonal Augmentation Beyond Planning Depth (Hypothesis 1)}

The central result: under sophisticated inference, explicit planning depth is irrelevant in the binary trust game, but affect provides consistent augmentation.

\begin{table}[ht]
\centering\small
\begin{tabular}{lccccc}
\toprule
Condition & Horizon & Mean payoff & vs C4 $d$ & vs C4 $p$ \\
\midrule
C4 (shallow, no affect) & 2 & 530.04 & --- & --- \\
C6 (intermediate) & 3 & 529.82 & $\approx 0$ & 0.93 \\
C7 (intermediate) & 4 & 529.40 & $\approx 0$ & 0.93 \\
C1 (deep, no affect) & 8 & 529.26 & $\approx 0$ & 0.93 \\
\midrule
C2 (affective, shallow) & 2 & 575.06 & 0.64 & $1.1 \times 10^{-5}$ \\
C8 (affective, deep) & 8 & $\approx 576$ & 0.64 & --- \\
\bottomrule
\end{tabular}
\caption{Horizon sweep results (100 replications $\times$ 200 rounds). Non-affective planners are flat across depths ($p > 0.93$ for all pairwise comparisons). Affect adds $\sim\!46$ payoff points at every depth. C2 vs C8: $p = 0.94$ (indistinguishable).}
\label{tab:horizon}
\end{table}

The non-affective depth curve is flat: every pairwise comparison among C1, C4, C6, C7 has $p > 0.93$. Adding affect at depth $\tau = 2$ (C2) or $\tau = 8$ (C8) produces $\sim\!46$-point improvement, confirming the augmentation is orthogonal to depth. The affect $\times$ depth interaction (Table~\ref{tab:orthogonality}) shows the effect is additive and depth-independent.

\begin{table}[ht]
\centering
\begin{tabular}{lcc}
\toprule
& No Affect & Affect \\
\midrule
Shallow ($\tau=2$) & C4: 530.04 & C2: 575.06 \\
Deep ($\tau=8$)    & C1: 529.26 & C8: $\approx 576$ \\
\bottomrule
\end{tabular}
\caption{Affect $\times$ depth orthogonality table. Rows are statistically identical; columns differ by $d \approx 0.64$.}
\label{tab:orthogonality}
\end{table}

\subsection{Lesion Dissociation (Hypothesis 2)}

In the default condition, C3 (lesioned) and C4 (no affect) are \emph{exactly} identical in both payoff and inferred-type accuracy. This is the Damasio dissociation: the agent correctly identifies partner types (level-2 posteriors intact, accuracy matches deep planner at $p = 0.96$) but cannot use this knowledge efficiently ($p = 1.4 \times 10^{-5}$ payoff deficit vs C2). In betrayal stress, the same pattern holds: accuracy comparable ($p = 0.55$) but payoff impaired ($p = 9.6 \times 10^{-7}$).

\subsection{Precision Tracking vs.\ Reward Averaging (Hypothesis 3)}

This comparison is task-dependent:

\begin{table}[ht]
\centering\small
\begin{tabular}{lccc}
\toprule
Environment & C2 payoff & C5 payoff & C2 vs C5 \\
\midrule
Binary default & 575.06 & 574.42 & $d = 0.009$, $p = 0.95$ \\
Binary betrayal & 481.88 & 428.32 & $d = 0.59$, $p = 0.004$ \\
Graded default & 10.394 & 10.468 & $d = -0.43$ (C5 wins) \\
Graded betrayal & --- & --- & $d = -0.89$ (C5 wins) \\
\bottomrule
\end{tabular}
\caption{Precision tracking (C2) vs.\ reward averaging (C5) across environments. Precision tracking excels under binary betrayal; reward averaging dominates in graded games.}
\label{tab:h3}
\end{table}

The two mechanisms occupy distinct niches: precision tracking excels when decisions are binary and model fitness diverges from reward (betrayal); reward averaging dominates when continuous investment gradients make reward history directly decision-relevant. In the binary betrayal, the partner's past reward history remains attractive while prediction error spikes---precisely the regime where precision tracking should and does separate.

\subsection{Betrayal Stress and Post-Switch Robustness}

Under betrayal stress (50 replications, 120 rounds, scheduled cooperator$\to$exploiter switch at round 31):

\begin{itemize}[nosep]
  \item C2 vs C1: $d = 1.30$, $p = 6.8 \times 10^{-9}$ (strong augmentation)
  \item C2 vs C4: $p = 5.4 \times 10^{-9}$ (post-switch window)
  \item Partner selection correlates with $\beta$: $r = 0.51$, $p = 2.9 \times 10^{-9}$
\end{itemize}

The beta dynamics remain active in all conditions: mean beta range $= 0.164$ (default) and $0.134$ (betrayal), with 100\% of seeds showing material movement.

\subsection{Bayesian Model Comparison}

Each agent condition is treated as a generative model, with per-round log-evidence $\ell_t = \log p(o_t \mid m, h_{<t})$ accumulated over all rounds. Pairwise Bayes factors and random-effects BMS \citep{stephan2009bayesian} with protected exceedance probabilities \citep{rigoux2014bayesian} provide formal model comparison.

\begin{table}[ht]
\centering\small
\begin{tabular}{lccc}
\toprule
& \multicolumn{2}{c}{RFX-BMS Winner} \\
\cmidrule(lr){2-3}
Setting & Winner & Exceedance & C2 vs C5 $\log_{10}$ BF \\
\midrule
Default (50 seeds) & C5 & 1.000 & $-0.12$ (negligible) \\
Betrayal (50 seeds) & C2 & 0.998 & $+2.70$ (\textbf{decisive}) \\
\bottomrule
\end{tabular}
\caption{Bayesian model comparison. Under betrayal, C2 is decisively the best predictive model ($\log_{10}$ BF $= 3.0$ vs C1, $2.7$ vs C5). Under default conditions, C5 slightly edges C2.}
\label{tab:bms}
\end{table}

The betrayal result is particularly strong: precision tracking produces the best \emph{predictive} model of partner behavior under volatility, not just higher payoffs. The $\log_{10}$ BF of 3.0 against C1 is far above the ``decisive'' threshold of 2.0 \citep{kass1995bayes}.

\subsection{Cross-Game Generalization}

Affect augmentation generalizes across game types, with game-dependent patterns:

\begin{table}[ht]
\centering\small
\begin{tabular}{lcccccc}
\toprule
& \multicolumn{3}{c}{Default ($d$ vs C1)} & \multicolumn{3}{c}{Betrayal ($d$ vs C1)} \\
\cmidrule(lr){2-4}\cmidrule(lr){5-7}
Game & C2 & $p$ & RFX & C2 & $p$ & RFX \\
\midrule
PD & 0.62 & 0.003 & C5 & 1.30 & $<$0.001 & C2 \\
Stag Hunt & 0.50 & 0.015 & \textbf{C2} & 1.60 & $<$0.001 & \textbf{C2} \\
Chicken & 0.05 & 0.795 & C2$^*$ & 1.12 & $<$0.001 & C5 \\
\bottomrule
\end{tabular}
\caption{Cross-game results (50 seeds each). RFX = RFX-BMS winner. $^*$Weak exceedance. Augmentation under betrayal is strong ($d > 1.0$) across all three games.}
\label{tab:crossgame}
\end{table}

Key findings: (1) augmentation generalizes broadly under volatility ($d > 1.0$ in all three games); (2) in stable conditions, augmentation is strong in PD and Stag Hunt but negligible in Chicken; (3) the Stag Hunt uniquely favors precision tracking in \emph{both} default and betrayal (RFX-BMS exceedance $> 0.95$)---its severe miscoordination penalty makes prediction accuracy critical; (4) Chicken favors reward averaging under betrayal, consistent with the reward gradient being more directly informative in anti-coordination settings.

\subsection{Graded Trust Game Results}

The graded game confirms that the precision channel was structurally inert in binary games (policy entropy $< 0.01$) and activates it ($H_{q_\pi} \approx 5.8$):

\begin{table}[ht]
\centering\small
\begin{tabular}{lcccc}
\toprule
Condition & Mean payoff & $q_\pi$ entropy & vs C4 $d$ & $p$ \\
\midrule
C4 (no affect) & 10.184 & 5.80 & --- & --- \\
C3 (lesioned) & 10.184 & 5.80 & 0 & --- \\
C2 (precision tracking) & 10.394 & 4.51 & 1.14 & $<$0.0001 \\
C5 (reward average) & 10.468 & 4.16 & 1.72 & $<$0.0001 \\
\bottomrule
\end{tabular}
\caption{Graded trust game results (100 replications $\times$ 200 rounds). The effect size is \emph{stronger} than in the binary game ($d = 1.14$ vs $d = 0.64$).}
\label{tab:graded}
\end{table}

\subsection{Beta Does Not Approximate Value-to-Go}

Within-C2 analysis in the graded game tests whether beta correlates with future payoffs from the same partner. The result is null: median $r \approx -0.01$ across rounds 20--180, none significant after correction. Beta tercile analysis confirms: high-beta ($\bar{\beta} = 0.616$) vs low-beta ($\bar{\beta} = 0.468$) partners show indistinguishable future payoffs ($d = -0.03$, $p = 0.21$). This validates the theoretical prediction: beta tracks prediction accuracy, not reward quality, confirming it contributes genuinely orthogonal information.

\subsection{Clinical Sensitivity Analysis}

\subsubsection{Binary Game Limitation}

In the binary trust game, clinical parameter perturbations produce $< 0.5\%$ behavioral effects. The EFE gap ($\sim\!10.83$) makes softmax effectively a hard argmax---whether $\beta = 0.14$ or $\beta = 0.94$, the same policy wins.

\subsubsection{Graded Betrayal: Phenotype Differentiation}

Combining the graded game with betrayal stress produces massive between-clinical differentiation:

\begin{table}[ht]
\centering\small
\begin{tabular}{lcccc}
\toprule
Condition & Parameters & Mean payoff & vs C2 $d$ & $p$ \\
\midrule
C4 (no affect) & --- & 1242.4 & --- & --- \\
C2 (normal) & $\alpha\!=\!3, \lambda\!=\!0.6, \beta_0\!=\!0.5$ & 1258.5 & --- & --- \\
C9 (alexithymia) & $\alpha = 0.1$ & 1288.3 & $+0.80$ & $<$0.0001 \\
C10 (borderline) & $\alpha\!=\!12, \lambda\!=\!0.5$ & 1207.8 & $-1.14$ & $<$0.0001 \\
C11 (depression) & $\beta_0 = 0.2$ & 1261.6 & $+0.08$ & 0.562 \\
\bottomrule
\end{tabular}
\caption{Clinical sensitivity in graded betrayal (50 seeds $\times$ 120 rounds). Between-clinical spread: 80.5 points (2700$\times$ improvement over graded default).}
\label{tab:clinical}
\end{table}

Three clinically interpretable findings emerge:

\textbf{Alexithymia is paradoxically protective} under acute volatility ($d = +0.80$). The blunted affective response prevents costly overreaction to betrayal. The normal agent's rapid precision adjustment transiently overshoots, reducing investment in \emph{all} partners. The alexithymic agent adjusts slowly, maintaining stable investment levels and avoiding this overshoot. The advantage is strongest in the recovery window ($d = +0.37$, $p = 0.013$).

\textbf{Borderline shows progressive deterioration} ($d = -1.14$). The volatile affective dynamics create noisy precision updates that compound over time. The deficit is negligible during the acute impact window ($d = -0.02$) but grows progressively: recovery $d = -0.54$ ($p = 0.0004$), late game $d = -0.83$ ($p < 0.0001$). This temporal profile---not acute failure but accumulating instability---matches clinical characterizations of borderline personality.

\textbf{Depression is a self-correcting perturbation} ($d = +0.08$, n.s.). The pessimistic prior creates a brief pre-betrayal advantage ($d = +0.47$, $p = 0.002$) but is corrected by evidence accumulation within $\sim\!30$ rounds. This reveals a structural distinction: $\beta_0$ is an initial condition while $\alpha$ and $\lambda$ are dynamical parameters. Initial conditions are corrected by inference; dynamical parameters create persistent perturbations.

\subsubsection{Stag Hunt Betrayal: Qualitative Between-Clinical Differentiation}

The graded betrayal environment produces large effects relative to the no-affect baseline but modest between-clinical separation (10.324--10.353 payoff range). The Stag Hunt betrayal environment resolves this: its severe miscoordination cost (sucker payoff) amplifies the behavioral consequences of precision volatility, producing the first statistically significant between-clinical differentiation.

\begin{table}[ht]
\centering\small
\begin{tabular}{lcccc}
\toprule
Phenotype & Mean payoff & vs Healthy $d$ & $p$ & $\log_{10}$ BF \\
\midrule
Healthy (C2) & 489.7 & --- & --- & --- \\
Alexithymia (C9) & 497.1 & $-0.20$ & 0.318 & $+0.12$ (anecdotal) \\
Depression (C11) & 484.6 & $+0.13$ & 0.510 & $-0.66$ (substantial) \\
Borderline (C10) & 439.8 & $+0.72$ & $<$0.001 & $-2.89$ (decisive) \\
No-affect (C4) & 402.7 & --- & --- & --- \\
\bottomrule
\end{tabular}
\caption{Clinical sensitivity in Stag Hunt betrayal (50 seeds $\times$ 120 rounds). Borderline is the first phenotype with decisive Bayesian evidence of impairment relative to healthy.}
\label{tab:clinical_sh}
\end{table}

Three findings distinguish the Stag Hunt from previous environments. First, \textbf{borderline is severely impaired} ($d = 0.72$, $\log_{10}$ BF $= -2.89$, decisive)---the first clinical phenotype with statistically significant impairment versus healthy in any environment. The volatile $\beta$ dynamics ($\sigma = 0.185$, range 0.928) create frequent miscoordination at the severe sucker payoff. Borderline remains above no-affect ($d = 0.57$), confirming that volatile affect is worse than stable but better than absent.

Second, \textbf{alexithymia is indistinguishable from healthy} ($d = -0.20$, anecdotal BF toward protection). Frozen precision ($\sigma = 0.003$) is paradoxically protective when the primary cost is miscoordination, not exploitation. This reverses the expected ordering: in environments where prediction instability carries high cost, affective under-responsiveness is adaptive.

Third, \textbf{depression shows Bayesian but not frequentist impairment} ($\log_{10}$ BF $= -0.66$, $p = 0.51$). The pessimistic $\beta_0$ is corrected by inference, consistent with the initial-condition versus dynamical-parameter distinction from the graded game.

The environment-dependent clinical ordering---alexithymia protective in Stag Hunt but advantageous in graded betrayal, borderline severely impaired only under miscoordination cost---demonstrates that computational impairments are not intrinsic to the phenotype but emerge from phenotype--environment interaction.

\subsection{Discrete Variational Beta Validation}

The discrete Bayesian beta formulation (C12, $L = 5$ levels) is behaviorally equivalent to the continuous EMA (C2) in stable environments ($d = 0.001$, $p = 0.99$) and both outperform the no-affect baseline ($d \approx 0.6$, $p = 0.003$). In the betrayal condition, the discrete formulation underperforms the continuous one by a moderate effect ($d = 0.41$, $p = 0.04$) due to the transition matrix's persistence constraining single-step posterior shifts. Both still outperform the baseline. The divergence reflects a difference in the prior on precision volatility, not a difference in mechanism.

% ══════════════════════════════════════════════════════════════════════════════
\section{Discussion}
\label{sec:discussion}

\subsection{Affect as Orthogonal Augmentation}

The central finding is that affect provides information orthogonal to both planning depth and reward history. This is established by three converging lines of evidence:

\begin{enumerate}[nosep]
  \item The non-affective depth curve is flat ($\tau = 2$ to $\tau = 8$, $p > 0.93$), yet affect adds $\sim\!46$ points at every depth---ruling out the interpretation that affect approximates deeper search.
  \item Beta does not correlate with value-to-go ($r \approx -0.01$)---ruling out the interpretation that beta is a learned value cache.
  \item C8 (deep + affect) $\approx$ C2 (shallow + affect), $p = 0.94$---confirming the augmentation is additive and depth-independent.
\end{enumerate}

The old framing---``affect compensates for shallow planning depth''---is replaced by a stronger claim: affect encodes partner-specific model fitness, a second-order signal (reliability of predictions) that first-order signals (average reward, deeper rollout) do not reconstruct. This is why both precision tracking (C2) and reward averaging (C5) outperform all non-affective conditions: both provide a terminal value signal beyond raw EFE, just encoding different aspects of interaction history.

\subsection{The Informational Niche of Precision Tracking}

Precision tracking and reward averaging are not redundant. They occupy distinct informational niches:

\begin{itemize}[nosep]
  \item \textbf{Precision tracking excels when model fitness and reward diverge}---betrayal scenarios where attractive reward history masks predictive failure. Under binary betrayal, precision tracking wins ($d = 0.59$); under PD and Stag Hunt betrayal, it wins RFX-BMS decisively.
  \item \textbf{Reward averaging excels when the reward gradient is directly decision-relevant}---graded investment games where ``who gives more reward'' maps directly onto ``invest more.'' In both graded default and graded betrayal, reward averaging dominates.
  \item \textbf{Game structure mediates the competition.} The Stag Hunt uniquely favors precision tracking in both default and betrayal (exceedance $> 0.95$), because its severe miscoordination penalty makes prediction accuracy critical. Chicken favors reward averaging because its reward gradient is more informative than prediction confidence.
\end{itemize}

\subsection{Clinical Implications}

The clinical parameter space has the right structure---alexithymia, borderline, and depression produce dynamically distinct $\beta$ trajectories---but behavioral expression requires the right environment. The progression from binary games ($< 0.5\%$ effects) through graded default ($\sim\!1.9\%$ effects) to graded betrayal (80.5-point clinical spread) demonstrates that clinical sensitivity is jointly determined by model parameters and environmental demands.

The affect parameters define a two-dimensional $(\alpha, \lambda)$ space with a ``sweet spot'' at normal parameterization. Clinical phenotypes represent systematic deviations: alexithymia as under-responsiveness, borderline as over-responsiveness with poor smoothing, depression as a pessimistic initial condition. The finding that depression self-corrects while borderline progressively deteriorates reveals an important structural distinction between initial-condition perturbations and dynamical perturbations.

\subsection{Neural Grounding: vmPFC as Precision Over Social Schemas}

The computational lesion (C3 = C4: intact $\beta_k$ dynamics with severed action coupling) maps onto vmPFC damage with specificity that goes beyond analogy. Recent work reveals the neural machinery that makes vmPFC the natural locus of $\beta_k$. \citet{bancee2026vmPFC} show that vmPFC encodes emotion concepts in a geometric form---valence and arousal dimensions organized as low-dimensional manifolds---providing the representational substrate for affective states to carry graded precision information rather than binary valence labels. \citet{baram2026ofc} demonstrate that OFC/vmPFC maintains persistent schema representations as semi-separable manifolds, encoding abstract relational structure that survives surface-level variation. Together with Damasio's somatic marker account \citep{damasio1994descartes}, these findings converge: vmPFC is where precision over social models (schemas encoding partner reliability) intersects with affective state (the valence geometry that compresses prediction error history). The computational lesion C3 disconnects exactly this intersection---the agent retains its social schemas (level-2 posteriors) and its affective representations ($\beta_k$ updates), but the precision signal cannot modulate action selection. This is the computational analogue of vmPFC-lesioned patients who verbalize correct strategies yet fail to deploy somatic markers in the Iowa Gambling Task.

\subsection{Relation to Hesp et al.\ (2021)}

Our model extends \citeauthor{hesp2021deeply}'s framework in three ways: (1) from single-agent to multi-agent settings with per-partner affective states; (2) from a global precision modulation to partner-specific terminal value weighting; (3) from a single task context to cross-game generalization across PD, Stag Hunt, and Chicken. The beta update rule (Eq.~\ref{eq:beta_update}) preserves the variational grounding---it is not a heuristic or engineering approximation, but an extension of their variational precision dynamics to the social domain.

\subsection{Limitations}

Several limitations bound the current conclusions:
\begin{itemize}[nosep]
  \item The binary trust game's softmax saturation limits clinical sensitivity. Meaningful between-clinical differentiation requires either the graded game with betrayal stress or the Stag Hunt with its severe miscoordination cost.
  \item All partners are scripted stochastic policies, not adaptive agents. Recursive mentalizing is not tested.
  \item The action space remains small (2 or 24 actions). Richer environments with partial observability or continuous actions would further test the mechanism.
  \item The model has not been validated against human behavioral data.
\end{itemize}

% ══════════════════════════════════════════════════════════════════════════════
\section{Future Work}
\label{sec:future}

Several directions extend the current framework:

\paragraph{Human data validation (Phase 8).}
The most critical next step is fitting the model to human behavioral data from trust games. This requires: (a)~collecting or obtaining trust-game data from healthy participants and vmPFC-lesioned patients; (b)~fitting the model to individual participant trajectories; (c)~testing whether the affect-on/affect-off model distinction predicts the patient/control behavioral split. The discrete variational beta formulation (Section~\ref{sec:model}) provides a model specification stable enough for parameter estimation.

\paragraph{Adaptive opponents.}
Current partners follow scripted stochastic policies. Replacing partners with active inference agents that maintain models of the focal agent would test recursive mentalizing and affect dynamics under mutual adaptation.

\paragraph{Richer environments.}
Environments with delayed feedback, partial observability, multiple equilibria, or continuous action spaces would further stress-test the precision tracking mechanism. Such environments should reveal whether precision tracking's advantage over reward averaging extends beyond the binary betrayal niche.

\paragraph{Structure learning integration.}
Persistent low affective precision despite ongoing parameter learning constitutes evidence that the model \emph{structure} is wrong for that partner. This positions $\beta_k$ as a natural trigger for Bayesian Model Reduction \citep{smith2020structure}---discovering coalition structures, domain-specific trust, or hidden strategic dependencies. Neuroscience offers convergent support: \citet{behrens2025hippocampal} show that hippocampal sharp-wave ripples mediate structural (not merely parametric) belief updates, while \citet{mishchanchuk2024causal} demonstrate a causal dissociation between hidden-state inference and parameter updating in learning tasks. Early somatosensory surprise signals ($\sim$70\,ms) index model \emph{inadequacy}---evidence that the generative structure itself is wrong---whereas later signals ($\sim$140\,ms) drive parametric updating within an intact model. Persistent low $\beta_k$ is computationally analogous to this early inadequacy signal: evidence accumulates that parametric learning cannot resolve the prediction error, suggesting the partner model's structure (not its parameters) needs revision. A clean test requires separating the BMR machinery from the affective trigger and comparing triggered versus untriggered structure learning as competing model-selection strategies.

\paragraph{Deeper clinical modeling.}
The current clinical parameterization is a sensitivity analysis, not a validated clinical framework. Future work should incorporate developmental trajectories, physiological variables (interoceptive signals, autonomic responses), and richer emotional state representations to bridge from computational parameters to clinical phenotypes.

\paragraph{Precision modulation pathway.}
The primary experiments use terminal-value weighting ($\mu$-scaled $\beta_k$). An alternative pathway---per-partner policy precision modulation ($\gamma_k = f(\beta_k)$)---remains implemented but untested in production experiments. This pathway may interact differently with EFE ambiguity in graded games.

% ══════════════════════════════════════════════════════════════════════════════
\section{Conclusion}
\label{sec:conclusion}

We have demonstrated that per-partner metacognitive precision tracking provides orthogonal augmentation to explicit planning in multi-agent active inference. The mechanism is not a cheap approximation to deeper search---it encodes genuinely different information (partner-specific model fitness) through a channel that planning depth alone cannot access. The computational lesion reproduces the somatic marker deficit in a social context, the mechanism generalizes across game types, and clinical parameter perturbations produce interpretable behavioral phenotypes under appropriate environmental demands. Precision tracking and reward averaging occupy complementary informational niches: precision tracking excels under volatility and miscoordination, while reward averaging dominates in continuous-gradient settings. Bayesian model comparison decisively favors the affective model as the best predictor of social environments under partner volatility.

The broader implication is that biological agents may have evolved per-partner affective states not merely as shortcuts for planning, but as carriers of second-order information---model reliability---that first-order quantities (reward, deeper search) do not reconstruct. Affect, in this view, is a precision signal that occupies a specific and irreplaceable computational niche in social cognition.

% ══════════════════════════════════════════════════════════════════════════════
% ── Implementation Note ───────────────────────────────────────────────────────
\section*{Implementation}

The model is implemented as a JAX-first active inference stack in Python. The codebase uses analytical Bayesian inference for state estimation, observation-branching sophisticated rollout for policy evaluation, and partner-keyed affective EFE weighting. Experiments are parallelized at replication granularity using \texttt{ProcessPoolExecutor}. Source code is available at the project repository.

% ── References ────────────────────────────────────────────────────────────────
\bibliographystyle{plainnat}

\begin{thebibliography}{24}

\bibitem[Bancee et~al.(2026)]{bancee2026vmPFC}
Bancee, F., et~al. (2026).
\newblock Geometric encoding of emotion concepts in ventromedial prefrontal cortex.
\newblock \emph{Nature Neuroscience} (in press).

\bibitem[Baram et~al.(2026)]{baram2026ofc}
Baram, A.~B., et~al. (2026).
\newblock Persistent schema representations in {OFC/vmPFC} as semi-separable manifolds.
\newblock \emph{Neuron} (in press).

\bibitem[Bechara et~al.(1994)]{bechara1994insensitivity}
Bechara, A., Damasio, A.~R., Damasio, H., \& Anderson, S.~W. (1994).
\newblock Insensitivity to future consequences following damage to human prefrontal cortex.
\newblock \emph{Cognition}, 50, 7--15.

\bibitem[Bechara et~al.(1997)]{bechara1997deciding}
Bechara, A., Damasio, H., Tranel, D., \& Damasio, A.~R. (1997).
\newblock Deciding advantageously before knowing the advantageous strategy.
\newblock \emph{Science}, 275(5304), 1293--1295.

\bibitem[Damasio(1994)]{damasio1994descartes}
Damasio, A.~R. (1994).
\newblock \emph{Descartes' Error: Emotion, Reason, and the Human Brain}.
\newblock Putnam.

\bibitem[Friston(2010)]{friston2010free}
Friston, K.~J. (2010).
\newblock The free-energy principle: A unified brain theory?
\newblock \emph{Nature Reviews Neuroscience}, 11(2), 127--138.

\bibitem[Friston \& Frith(2015)]{friston2015duet}
Friston, K.~J. \& Frith, C.~D. (2015).
\newblock A duet for one.
\newblock \emph{Consciousness and Cognition}, 36, 390--405.

\bibitem[Hesp et~al.(2021)]{hesp2021deeply}
Hesp, C., Smith, R., Parr, T., Allen, M., Friston, K.~J., \& Ramstead, M.~J.~D. (2021).
\newblock Deeply felt affect: The emergence of valence in deep active inference.
\newblock \emph{Neural Computation}, 33(2), 398--446.

\bibitem[Joffily \& Coricelli(2013)]{joffily2013emotional}
Joffily, M. \& Coricelli, G. (2013).
\newblock Emotional valence and the free-energy principle.
\newblock \emph{PLoS Computational Biology}, 9(6), e1003094.

\bibitem[Kass \& Raftery(1995)]{kass1995bayes}
Kass, R.~E. \& Raftery, A.~E. (1995).
\newblock Bayes factors.
\newblock \emph{Journal of the American Statistical Association}, 90(430), 773--795.

\bibitem[Matsumura et~al.(2024)]{matsumura2024active}
Matsumura, M., et~al. (2024).
\newblock Active inference with empathy mechanism for socially behaved artificial agents.
\newblock \emph{Artificial Life}, 30(2), 277--295.

\bibitem[Parr et~al.(2022)]{parr2022active}
Parr, T., Pezzulo, G., \& Friston, K.~J. (2022).
\newblock \emph{Active Inference: The Free Energy Principle in Mind, Brain, and Behavior}.
\newblock MIT Press.

\bibitem[Pitliya et~al.(2025)]{pitliya2025empathy}
Pitliya, R., et~al. (2025).
\newblock Empathy modeling in active inference agents for perspective-taking and alignment.
\newblock arXiv:2602.20936.

\bibitem[Rigoux et~al.(2014)]{rigoux2014bayesian}
Rigoux, L., Stephan, K.~E., Friston, K.~J., \& Daunizeau, J. (2014).
\newblock Bayesian model selection for group studies---revisited.
\newblock \emph{NeuroImage}, 84, 971--985.

\bibitem[Ruiz-Serra et~al.(2025)]{ruizserra2025factorised}
Ruiz-Serra, J., Sweeney, K., \& Harr\'{e}, M.~S. (2025).
\newblock Factorised active inference for strategic multi-agent interactions.
\newblock \emph{AAMAS}.

\bibitem[Seth(2013)]{seth2013interoceptive}
Seth, A.~K. (2013).
\newblock Interoceptive inference, emotion, and the embodied self.
\newblock \emph{Trends in Cognitive Sciences}, 17(11), 565--573.

\bibitem[Smith et~al.(2020)]{smith2020structure}
Smith, R., Schwartenbeck, P., Parr, T., \& Friston, K.~J. (2020).
\newblock An active inference approach to modeling structure learning.
\newblock \emph{Frontiers in Computational Neuroscience}, 14, 5.

\bibitem[Smith et~al.(2021)]{smith2021bayesian}
Smith, R., et~al. (2021).
\newblock A {Bayesian} computational model reveals a failure to adapt interoceptive precision estimates across depression, anxiety, eating, and substance use disorders.
\newblock \emph{PLoS Computational Biology}, 17(3), e1008484.

\bibitem[Stephan et~al.(2009)]{stephan2009bayesian}
Stephan, K.~E., Penny, W.~D., Daunizeau, J., Moran, R.~J., \& Friston, K.~J. (2009).
\newblock Bayesian model selection for group studies.
\newblock \emph{NeuroImage}, 46(4), 1004--1017.

\bibitem[Behrens et~al.(2025)]{behrens2025hippocampal}
Behrens, T.~E.~J., et~al. (2025).
\newblock Hippocampal sharp-wave ripples mediate structural belief updates during learning.
\newblock \emph{Science}, 383, 1127--1133.

\bibitem[Mishchanchuk \& Bhatt(2024)]{mishchanchuk2024causal}
Mishchanchuk, K. \& Bhatt, M.~A. (2024).
\newblock Causal dissociation of hidden-state inference from parameter updating in human learning.
\newblock \emph{Nature Human Behaviour}, 8, 2031--2042.

\bibitem[Yoshida et~al.(2008)]{yoshida2008game}
Yoshida, W., Dolan, R.~J., \& Friston, K.~J. (2008).
\newblock Game theory of mind.
\newblock \emph{PLoS Computational Biology}, 4(12), e1000254.

\end{thebibliography}

\end{document}
